% ============================================================
% Chapter 3: Project Management
%   Section: Planning
%     Subsection: Phases and Tasks
% Included from chapters/03-project-management/02-planning/planning.tex
% ============================================================

\subsection{Phases and Tasks}
\label{subsec:phases}

\begin{table}[H]
\centering
\small
\caption{Phases, dates and planned hours (docs/roadmap.md, rebalanced
on 30 September 2026).}
\label{tab:phases}
\begin{tabularx}{\linewidth}{@{}lXlr@{}}
\toprule
Phase & Content & Dates (2026) & Hours \\
\midrule
0 & Foundations: scraping core, enrichment, fact-checking v1, job
    API and frontend v1 & Jun 1 -- Jul 31 & 300 \\
-- & Vacation pause & Aug 1 -- 31 & 0 \\
1 & Pipeline completion: stabilisation, testing and bugfixing
    & Sep 1 -- 28 & 140 \\
2 & Development: scraping strategies (Sprint 3) and fact-checking
    improvements (Sprint 4) & Sep 29 -- Oct 19 & 120 \\
3 & Neo4j (done on Sep 29) and infrastructure deployment
    & Oct 20 -- Nov 9 & 50 \\
4 & Evaluation: labelling, harness and tuning & Sep 29 -- Nov 16 & 90 \\
5 & AI model testing: writing (Sprint 6) and verification
    (Sprint 7) models & Nov 10 -- Dec 10 & 80 \\
6 & Interface: reader view & Dec 11 -- 17 & 35 \\
-- & Testing and QA, transversal & Sep 29 -- Dec 24 & 50 \\
7 & Thesis report & Dec 25 -- 31 & 35 \\
\midrule
  & \textbf{Total} & & \textbf{900} \\
\bottomrule
\end{tabularx}
\end{table}

\paragraph{Phase 0 --- Foundations.} Four sprints built the pipeline
end to end: the repository, the API skeleton and the declarative source
catalogue with RSS discovery and trafilatura extraction (Sprint~0.1);
the enrichment stack --- keywords, entities, claims, topics, sentiment,
quality and embeddings (0.2); the first fact-checking pipeline, with
admission, claim selection, SearXNG and vector retrieval, ranking, the
language-model verifier and the confidence scorer (0.3); and the job
API with polling, the text corrector and the first frontend (0.4).

\paragraph{Phase 1 --- Pipeline completion.} Sprint~1 stabilised the
core: pipeline edge cases, progress events for every stage, the
correctness of the analysis cache and the shared test fakes. Sprint~2
was testing and bug-fixing: a full regression pass, concurrency tests,
a load test of the synchronous endpoint and hardened handling of
language-model failures. An audit on 21~September found and fixed a
re-analysed URL being rejected as a duplicate of itself, a committed
database password and the loss of job history on restart, and made the
claims run concurrently.

\paragraph{Phase 2 --- Development.} Sprint~3 built the scraping
strategies: an ingestion path, the BeautifulSoup and headless-browser
steps of the extraction cascade, topic-page discovery, a source probe,
and the catalogue's growth from 12 to 36 sources. Sprint~4 improved
fact-checking --- the search fix with a ``search unavailable'' state
and a keyless second route, alternative retrieval strategies, and
regression tests for ranking and confidence --- and added, outside the
original plan, discovery by topic group, the choice of up to twenty
candidates by a person or the AI, the newsroom flow in the frontend,
the mission screen, the candidate ranking and the topic emphasis
(Section~\ref{sec:deviations}).

\paragraph{Phase 3 --- Neo4j and infrastructure.} The graph schema, its
write path from every analysis, a backfill from the lake and the
labelled facts, a read use-case and its live tests (done on
29~September); then the hosting decision, a production compose file,
SearXNG for production, logging and monitoring, and the Google Cloud
infrastructure as Terraform, checked against a local emulator (done on
1--2~October, except the first real deployment).

\paragraph{Phase 4 --- Evaluation.} The custom validation set (daily
labelling batches, 150 facts as the target, a blind 20\% re-label and
its agreement), the evaluation harness (record format, resumable runs,
metrics with bootstrap intervals, retrieval evaluation, stage
attribution and journalist-effort metrics), a pilot on X-Fact, and the
tuning of thresholds against the labels.

\paragraph{Phase 5 --- AI model testing.} Sprint~6 benchmarks a local
and a hosted model as the editor behind the text corrector; Sprint~7
benchmarks language models as verifiers on the labelled sets, compares
their verdicts with the human labels, re-tunes the confidence weights
and selects the default model.

\paragraph{Phase 6 --- Interface.} A read-only feed of the published
articles by topic, an article view with each claim's verdict, evidence
and the reason each source was trusted, built on the existing design
system, and public access to that view only.

\paragraph{Testing and quality assurance.} Transversal: unit,
concurrency and graph tests, continuous integration, a full-pipeline
integration test, tests of the evaluation harness itself, a stress test
of the guardrails against a model that bluffs, infrastructure smoke
tests, load tests, and a final regression and bug-fix week.

\paragraph{Phase 7 --- Report.} Writing this document on a system that
no longer changes.

\paragraph{The unfunded 18 hours.} On 5~October an 18-hour block of
experiments on the stages before the verdict --- discovery, selection,
admission, the topic classifier and time per stage --- was added to
Phase~4 without hours to pay for it. It is reported as such: the plan
stands at 918 hours, an overrun of 2\%, unless the writing-model
benchmark of Sprint~6 (20 hours) is not run. The thesis priorities of
14~October (Section~\ref{sec:deviations}) moved that benchmark to
``only if time is left'' and moved ten of the eighteen hours --- the
discovery numbers and the blind selection rounds --- into the work
needed for RQ1 and RQ3, which is how the overrun is expected to be
absorbed.
